% Generated draft: requires mathematical and source review.
\documentclass[11pt]{article}
\usepackage[margin=25mm]{geometry}
\usepackage{fontspec}
\setmainfont{Latin Modern Roman}
\usepackage{amsmath,amssymb,mathtools,unicode-math}
\setmathfont{Latin Modern Math}
\usepackage{enumitem,needspace,xurl,hyperref}
\setlength{\parindent}{0pt}
\setlength{\parskip}{5pt}
\setlength{\emergencystretch}{3em}
\newcommand{\N}{\mathbb{N}}
\newcommand{\Z}{\mathbb{Z}}
\newcommand{\Q}{\mathbb{Q}}
\newcommand{\R}{\mathbb{R}}
\newcommand{\C}{\mathbb{C}}
\newcommand{\F}{\mathbb{F}}
\newcommand{\A}{\mathbb{A}}
\newcommand{\PP}{\mathbb{P}}
\newcommand{\E}{\mathbb{E}}
\newcommand{\eps}{\varepsilon}
\newcommand{\dd}{\,\mathrm{d}}
\begin{document}
\section*{UnsolvedMath ID: 700778}
\textbf{Source problem code:} catalogue entry
\subsection{Definitions and mathematical statement}
Let
\[
G=\langle x_1,\ldots,x_n\mid r_1=\cdots=r_m=1\rangle
\]
be a group given by a finite presentation. Write \(F_n=\langle x_1,\ldots,x_n\rangle\) for the free group on the displayed generators, and let \(N=\langle\!\langle r_1,\ldots,r_m\rangle\!\rangle_{F_n}\) be the normal closure of its relators, so that \(G=F_n/N\).

The source asks whether every such group possesses, ``in a natural way,'' a homomorphic image admitting a presentation
\[
H=\langle x_1,\ldots,x_n\mid s_1=\cdots=s_m=1\rangle
\]
with the same displayed numbers of generators and relators, satisfying each of the following requirements, considered separately:
\begin{enumerate}
\item[(a)] \(H\) is torsion-free: for every \(h\in H\) and every integer \(k\geq 1\), the equality \(h^k=1\) implies \(h=1\).
\item[(b)] The integral group ring \(\Z[H]\) embeds in a skew field. Here \(\Z[H]\) consists of finite formal sums \(\sum_{h\in H}a_hh\), with \(a_h\in\Z\), with multiplication induced by the group operation. A skew field is a nonzero, possibly noncommutative ring with identity in which every nonzero element is invertible; an embedding is understood as an injective identity-preserving ring homomorphism.
\end{enumerate}
If the repeated generator names mean that the quotient map sends each displayed generator of \(G\) to the corresponding displayed generator of \(H\), then, writing
\[
M=\langle\!\langle s_1,\ldots,s_m\rangle\!\rangle_{F_n},
\]
the quotient requirement is precisely \(N\subseteq M\). The source does not explicitly specify this compatibility, nor does it define ``in a natural way.'' Thus the displayed algebraic conditions do not fully specify the intended construction or its naturality requirements.

The source describes (b) as the stronger assertion and states that an affirmative answer would give an explicit method for finding displayed generators
\[
x_{i_1},\ldots,x_{i_{n-m}}
\]
that generate a free group in \(G\). It attributes the existence of such elements to a theorem of N. S. Romanovski\u{\i}. This accompanying claim requires qualification: no hypothesis \(n\geq m\) is stated, although the list has no meaning when \(n-m<0\). The supplied wording also does not explicitly say whether the selected generators must freely generate a subgroup of rank \(n-m\). The intended constructive content linking (b) to this method remains unspecified.
\subsection{Short English statement}
For every finitely presented group, can one naturally construct a quotient with a presentation having the same numbers of generators and relators, such that the quotient is torsion-free, or, more strongly as described by the source, its integral group ring embeds in a skew field? The meaning of naturality and the precise free-subgroup consequence require clarification.
\subsection{Sources}
\begin{itemize}
\item Local UnsolvedMath registry, supplied complete source record: permanent ID \(700778\), exact source problem code \texttt{catalogue entry}, title \emph{Kourovka Notebook Problem 11.86}. The registry labels the problem \texttt{open}; this is a supplied status label, not an independently verified mathematical assessment.
\item \emph{The Kourovka Notebook (Unsolved Problems in Group Theory)}, problem 11.86, proposed by V. A. Roman'kov in the 11th issue (1990). The supplied background also identifies the 21st issue edition (2026), \texttt{arXiv:1401.0300v48}. The registry describes the collection as edited by E. I. Khukhro and V. D. Mazurov. These bibliographic details are supplied attribution and have not been independently checked.
\item N. S. Romanovski\u{\i}, \emph{Algebra and Logic}, volume 16, no. 1 (1977), pp. 62--67. This is the citation supplied for the asserted existence of the indicated free subgroup; no article title or theorem statement is supplied.
\end{itemize}
\end{document}
